%! AP Statistics — Unit 4, Lesson 4.12 — Group Activity KEY
\documentclass[10pt]{article}
\usepackage{apstats-article}
\usepackage{apstats-key}

\begin{document}

\pageheader{Unit 4, Lesson 4.12}{Group Activity KEY: Geometric vs.\ Binomial}
\namepartnerperiod

\vspace{0.1in}

% ── TIER R ────────────────────────────────────────────────────────────────────
\begin{tcolorbox}[colback=goldbg, colframe=goldacc, boxrule=1pt, arc=2mm,
    title={\bfseries Tier R --- Key},
    fonttitle=\small]
\small

\textbf{Scenario 1:} Flipping until tails.
\begin{enumerate}[label=(\alph*), leftmargin=2em, itemsep=4pt]
  \item \ans{Geometric}
  \item \ans{$n$ is not fixed --- flipping continues until tails appears.}
  \item \ans{$X \sim \text{Geom}(0.5)$; $p = 0.5$}
\end{enumerate}

\textbf{Scenario 2:} Flipping 10 times.
\begin{enumerate}[label=(\alph*), leftmargin=2em, itemsep=4pt]
  \item \ans{Binomial}
  \item \ans{$n = 10$ is fixed in advance; $X$ counts the total number of tails.}
  \item \ans{$X \sim B(10, 0.5)$; $n = 10$, $p = 0.5$}
\end{enumerate}

\textbf{Scenario 3:} Number of first roll of a fair tetrahedral die.
\begin{enumerate}[label=(\alph*), leftmargin=2em, itemsep=4pt]
  \item \ans{Neither --- $X$ takes values $\{1, 2, 3, 4\}$ each with probability $1/4$. There is only one trial, so it is not binomial or geometric.}
  \item \ans{It is simply a discrete uniform distribution, not a repeated-trials distribution.}
\end{enumerate}
\end{tcolorbox}

\vspace{0.1in}

% ── TIER A ────────────────────────────────────────────────────────────────────
\begin{tcolorbox}[colback=sky, colframe=navy, boxrule=1pt, arc=2mm,
    title={\bfseries Tier A --- Key},
    fonttitle=\small]
\small

Archer, $p = 0.30$.

(a) BITS:
\begin{tabularx}{\linewidth}{@{} p{0.08\linewidth} X @{}}
B: & \ans{Binary --- each arrow either hits the bullseye (success) or misses (failure).} \\[4pt]
I: & \ans{Independent --- the outcome of one shot does not affect the next.} \\[4pt]
T: & \ans{Trials until first success --- shooting continues until the first bullseye.} \\[4pt]
S: & \ans{Same $p = 0.30$ on every shot.} \\[4pt]
\end{tabularx}
$X \sim \text{Geom}(0.30)$

(b) $P(X=5) = (0.70)^4(0.30) = (0.2401)(0.30) \approx \ans{0.0720}$

(c) $\mu_X = 1/0.30 \approx \ans{3.33}$

\ans{If the archer shoots arrows until hitting the bullseye many times, she will average about 3.33 arrows per bullseye.}

(d) $P(X \le 4) = \texttt{geometcdf(0.30, 4)} \approx \ans{0.7599}$

\ans{There is about a 76.0\% probability that the archer hits the bullseye within the first 4 shots.}
\end{tcolorbox}

\vspace{0.1in}

% ── TIER E ────────────────────────────────────────────────────────────────────
\begin{tcolorbox}[colback=greenbg, colframe=greenacc, boxrule=1pt, arc=2mm,
    title={\bfseries Tier E --- Key},
    fonttitle=\small]
\small

(a)
\renewcommand{\arraystretch}{1.8}
\begin{tabularx}{\linewidth}{@{} >{\bfseries}p{1.2cm} X X X X X @{}}
$k$ & 1 & 2 & 3 & 4 & 5 \\
\hline
$P(X=k)$ & \ans{0.3000} & \ans{0.2100} & \ans{0.1470} & \ans{0.1029} & \ans{0.0720} \\
\end{tabularx}

(b) \ans{The probabilities decrease as $k$ increases --- the distribution is strongly right-skewed with the highest probability at $k=1$. This contrasts sharply with a binomial distribution for the same $p$, which is roughly mound-shaped (bell-shaped) centered near $\mu_X = np$.}

(c) \ans{$P(X=5) \approx 0.0720$: the probability that the archer needs exactly 5 shots to get the first bullseye. $P(Y=5) = \binom{5}{5}(0.30)^5(0.70)^0 = (0.30)^5 \approx 0.0024$: the probability of getting exactly 5 bullseyes in 5 independent shots. They are different because $X=5$ (geometric) means 4 misses then a hit, while $Y=5$ (binomial) means all 5 shots are hits. Both use $p=0.30$ but count completely different events.}
\end{tcolorbox}

\end{document}
